\documentclass[11pt]{article}
\usepackage[margin=1in]{geometry}
\usepackage{amsmath,amssymb,microtype}
\usepackage[hidelinks]{hyperref}

\newcommand{\rad}{\operatorname{rad}}

\title{The hereditary-subalgebra semisimplicity question\\
in arXiv:0907.3103 was answered by arXiv:1206.3466}
\author{Literature-status note}
\date{13 August 2026}

\begin{document}
\maketitle

\begin{center}
\fbox{\parbox{0.91\textwidth}{\textbf{Status: literature already
answered.}  The later Theorem~4.2 proves the stronger exact formula for
the Jacobson radical of every hereditary subalgebra.}}
\end{center}

\section{The source question}

For the source, an HSA (hereditary subalgebra) $D$ of an operator algebra
$A$ is an approximately unital subalgebra satisfying $DAD\subset D$.
Proposition~2.5 of arXiv:0907.3103 proves that semiprimeness and topological
simplicity pass from $A$ to its HSA's.  Remark~2.6 immediately afterward,
on PDF page~4, states:
\begin{quote}
We do not know if semisimplicity passes to HSA's.
\end{quote}
Equivalently, if the Jacobson radical $J(A)$ vanishes, must $J(D)$ vanish
for every HSA $D\subset A$? \cite{ABS09}

\section{The exact later answer}

Theorem~4.2 of arXiv:1206.3466 (PDF page~7) states that, for every HSA
$D$ in every operator algebra $A$,
\[
 \boxed{J(D)=D\cap J(A).}
\]
It follows immediately that $J(A)=0$ implies $J(D)=0$; hence
semisimplicity does pass to HSA's \cite{ABR12}.

The later result is strictly stronger than the yes/no assertion.  Its proof
uses the spectral characterization of the radical.  One inclusion factors
$x\in J(D)$ as $x=dy$ by Cohen factorization and uses $DAD\subset D$ to
show $r(xa)=0$ for every $a\in A^1$.  The reverse inclusion uses the
preceding spectral-invariance proposition: quasi-invertibility of elements
of $D$ agrees whether computed in $D$ or in $A$.

The supporting paper introduces Theorem~4.2 with the explicit sentence
``The following answers a question posed in [3],'' and its reference [3] is
the source paper.  This is therefore an explicit later literature answer,
not an agent-generated theorem.  The other independent questions in
arXiv:0907.3103 are outside this packet's scope.

\begin{thebibliography}{9}
\bibitem{ABS09}
M.~Almus, D.~P. Blecher, and S.~Sharma,
\emph{Ideals and structure of operator algebras},
arXiv:0907.3103 (2009), Remark~2.6;
J. Operator Theory 67 (2012), 397--436.

\bibitem{ABR12}
M.~Almus, D.~P. Blecher, and C.~J. Read,
\emph{Ideals and hereditary subalgebras in operator algebras},
arXiv:1206.3466 (2012), Theorem~4.2;
Studia Math. 212 (2012), 65--93.
\end{thebibliography}

\end{document}
